\documentclass[12pt, a4paper]{article}

\usepackage[margin=1in]{geometry}
\usepackage{times}
\usepackage{graphicx}
\usepackage{booktabs}
\usepackage{amsmath}
\usepackage{algorithm}
\usepackage{algorithmic}
\usepackage{listings}
\usepackage{hyperref}
\usepackage{color}
\usepackage{authblk}

\definecolor{codegreen}{rgb}{0,0.6,0}
\definecolor{codegray}{rgb}{0.5,0.5,0.5}
\definecolor{codepurple}{rgb}{0.58,0,0.82}
\definecolor{backcolour}{rgb}{0.95,0.95,0.92}

\lstdefinestyle{mystyle}{
    backgroundcolor=\color{backcolour},   
    commentstyle=\color{codegreen},
    keywordstyle=\color{magenta},
    numberstyle=\tiny\color{codegray},
    stringstyle=\color{codepurple},
    basicstyle=\ttfamily\footnotesize,
    breakatwhitespace=false,         
    breaklines=true,                 
    captionpos=b,                    
    keepspaces=true,                 
    numbers=left,                    
    numbersep=5pt,                  
    showspaces=false,                
    showstringspaces=false,
    showtabs=false,                  
    tabsize=2
}
\lstset{style=mystyle}

\begin{document}

\title{Multi-threading vs.\ SIMD Performance Analysis in Smoothed Particle Hydrodynamics on ARM64 Architecture}

\author{Sujal Khatri}
\affil{Department of Computer Science, Fisk University}

\date{March 2026}

\maketitle

\begin{abstract}
Smoothed Particle Hydrodynamics (SPH) is a popular method for simulating fluids, but its $O(n^2)$ pairwise particle interactions make it very expensive to run in real time. In this project, we measured how much we can speed up a 2D SPH simulation on the Apple M1 processor using two different optimization approaches: multi-threading with OpenMP and SIMD vectorization with ARM NEON intrinsics. We built four versions of the same simulation---a scalar baseline, an OpenMP-parallelized version, a NEON SIMD version, and a combined version---and benchmarked all of them on the same 2,660-particle dam-break scenario. Our results show that OpenMP with 8 threads reaches 58.89 FPS (a 23.0$\times$ speedup over the unoptimized baseline), while single-threaded NEON SIMD reaches 49.69 FPS (19.4$\times$). The combined approach brings the total physics computation down from 6.169~ms to 3.185~ms per frame. We also found that restructuring particle data from Array of Structures (AoS) to Structure of Arrays (SoA) was critical for making SIMD work well, and that a single SIMD-optimized core can match roughly 85\% of the throughput of 8 threads running on unoptimized memory layouts.
\end{abstract}

\section{Introduction}

Smoothed Particle Hydrodynamics (SPH) is a mesh-free method for simulating fluids where instead of dividing space into a grid, the fluid is represented as a collection of moving particles \cite{muller2003particle}. Each particle carries properties like position, velocity, and density, and the simulation updates these properties every frame based on interactions with nearby particles. SPH works well for things like splashing water and free-surface flows because the particles naturally follow the fluid motion.

The big problem with SPH is performance. To compute each particle's density, we need to check its distance to every other particle and sum up kernel contributions from the ones that are close enough. In a simple (brute-force) implementation, this means an $O(n^2)$ loop---for 2,660 particles, that is about 7 million pair evaluations per frame. The force calculation has the same complexity. These two loops together dominate the runtime.

Modern processors like the Apple M1 offer two main ways to make this faster. First, the M1 has 8 CPU cores (4 fast ``Firestorm'' cores and 4 efficient ``Icestorm'' cores), so we can split the work across threads using OpenMP \cite{openmp2018specification}. Second, each core supports 128-bit ARM NEON SIMD instructions that can process 4 single-precision floats at once \cite{arm2020neon}. These two approaches are fundamentally different---one distributes work across cores (task parallelism), and the other makes each core do more work per instruction (data parallelism).

In this project, we wanted to understand how each of these approaches performs on its own, how they compare to each other, and what happens when we combine them. We also wanted to see how the choice of memory layout (AoS vs.\ SoA) affects the results. To do this, we built four versions of the same SPH simulation and benchmarked them all under identical conditions.

\subsection{Research Questions}
Our project targets four questions:
\begin{enumerate}
    \item How much speedup does multi-threading (OpenMP) provide on ARM64?
    \item What performance gains come from SIMD vectorization (ARM NEON)?
    \item How does switching from Array of Structures (AoS) to Structure of Arrays (SoA) affect cache efficiency?
    \item What are the main architectural bottlenecks for particle simulations on Apple Silicon?
\end{enumerate}

\section{Background}

\subsection{SPH Physics}
Our simulation follows the standard SPH formulation described by M\"uller et al.\ \cite{muller2003particle} and the Eurographics survey by Koschier et al.\ \cite{koschier2014sph}. Each frame, the simulation runs through three main steps.

\textbf{Step 1: Density.} For each particle $i$, we compute its local density by summing kernel contributions from all particles $j$ within the smoothing radius $h$:
\begin{equation}
\rho_i = \sum_{j=1}^{N} m_j \, W(|\mathbf{r}_i - \mathbf{r}_j|, h)
\end{equation}
We use the Poly6 kernel for density because it is smooth and only depends on the squared distance $r^2$ (so we can skip the expensive square root). Once we have density, we convert it to pressure using a simple equation of state:
\begin{equation}
P_i = k(\rho_i - \rho_0)
\end{equation}
where $k = 2000$ is the stiffness constant and $\rho_0 = 1.0$ is the rest density.

\textbf{Step 2: Forces.} Each particle accumulates three forces---pressure, viscosity, and gravity:
\begin{equation}
\mathbf{F}_i = -\sum_{j} m_j \left( \frac{P_i}{\rho_i^2} + \frac{P_j}{\rho_j^2} \right) \nabla W_{ij} + \mu \sum_{j} m_j \frac{\mathbf{v}_j - \mathbf{v}_i}{\rho_j} \nabla^2 W_{ij} + m_i \mathbf{g}
\end{equation}
The pressure gradient uses the Spiky kernel (which has a sharper peak than Poly6), and viscosity uses the Spiky Laplacian kernel. These kernels are taken directly from M\"uller et al.\ \cite{muller2003particle}.

\textbf{Step 3: Integration.} We update velocities and positions with semi-implicit Euler:
\begin{align}
\mathbf{v}_{t+\Delta t} &= \mathbf{v}_t + \frac{\mathbf{F}_t}{\rho_i} \Delta t \\
\mathbf{x}_{t+\Delta t} &= \mathbf{x}_t + \mathbf{v}_{t+\Delta t} \Delta t
\end{align}
Boundary collisions are handled by clamping positions and reflecting velocities with a damping factor of 0.5.

\subsection{Apple M1 Architecture}
The M1 chip \cite{apple2020m1} has a few properties that are relevant to our experiments:
\begin{itemize}
    \item \textbf{8 cores:} 4 Firestorm cores at 3.2 GHz and 4 Icestorm cores at 2.0 GHz.
    \item \textbf{128-bit NEON:} Each core can execute 4-wide float SIMD operations.
    \item \textbf{Unified memory:} CPU and GPU share the same physical memory, so there are no NUMA penalties when multiple cores access the same data.
    \item \textbf{Cache hierarchy:} 128 KB L1 data cache per core, up to 12 MB shared L2.
\end{itemize}
The unified memory is a double-edged sword---it avoids the coherency overhead of multi-socket systems, but all 8 cores share the same ~68 GB/s memory bandwidth, which can become a bottleneck under heavy parallel workloads.

\section{Methodology}

We built four versions of the simulation, each adding one optimization on top of the baseline. Every version uses the exact same physics parameters and the same 2,660-particle dam-break initial condition so we can compare them fairly.

\subsection{Phase 1: Baseline}
The baseline (\texttt{main\_improved.cpp}) is a straightforward single-threaded implementation using an Array of Structures (AoS) layout---each particle is a struct with position, velocity, force, density, and pressure fields. We compiled it with \texttt{-O0} (no compiler optimizations) to establish the true algorithmic cost without any help from the compiler.

This is intentionally slow. The point is to have a clean worst-case reference so we can measure how much each optimization actually helps.

\subsection{Phase 2: OpenMP Multi-threading}
The OpenMP version (\texttt{main\_omp.cpp}) keeps the same AoS layout and adds \texttt{\#pragma omp parallel for schedule(static)} to the outer \texttt{i} loop of both the density and force calculations. Since each iteration of the outer loop writes only to \texttt{particles[i]}, there are no write conflicts between threads---no locks or atomic operations needed.

We compiled this with \texttt{-O2} and linked against Homebrew's \texttt{libomp}. The simulation uses all 8 available threads. We also added timing instrumentation using \texttt{omp\_get\_wtime()} to measure the density and force loops separately.

\subsection{Phase 3: NEON SIMD}
The SIMD version (\texttt{main\_simd.cpp}) makes two changes: it restructures the data from AoS to SoA, and it vectorizes the inner \texttt{j} loop using ARM NEON intrinsics.

\textbf{SoA layout.} Instead of one array of \texttt{Particle} structs, we use separate arrays: \texttt{pos\_x[]}, \texttt{pos\_y[]}, \texttt{vel\_x[]}, etc. This means when the density loop reads positions, it gets 16 consecutive \texttt{pos\_x} values per 64-byte cache line instead of just 2 (because each AoS struct is 32 bytes, and only 8 bytes of that is position data).

\textbf{NEON vectorization.} The inner loop processes 4 particles at a time:
\begin{itemize}
    \item \texttt{vld1q\_f32}: loads 4 floats from a contiguous array
    \item \texttt{vsubq\_f32}, \texttt{vmulq\_f32}, \texttt{vmlaq\_f32}: arithmetic on 4-wide vectors
    \item \texttt{vcltq\_f32} + \texttt{vandq\_u32}: replaces the \texttt{if (r2 < h2)} branch with a branchless mask
    \item \texttt{vrecpeq\_f32} with Newton-Raphson: fast approximate reciprocal (avoids division)
    \item \texttt{vaddvq\_f32}: horizontal sum across all 4 lanes
\end{itemize}
A scalar remainder loop handles the last few particles when $n$ is not a multiple of 4. This version is compiled with \texttt{-O2 -mcpu=apple-m1} and runs on a single thread.

\subsection{Phase 4: Combined (SIMD + OpenMP)}
The combined version (\texttt{main\_combined.cpp}) takes the SoA + NEON code from Phase 3 and wraps the outer \texttt{i} loop with \texttt{\#pragma omp parallel for}. Each thread runs the NEON-vectorized inner loop on its assigned chunk of particles. Since threads only write to their own particle's output fields, there are no race conditions.

\subsection{Experimental Setup}
Table~\ref{tab:config} shows the fixed simulation parameters used across all benchmarks.

\begin{table}[h]
\centering
\caption{Simulation Configuration (identical across all phases)}
\label{tab:config}
\begin{tabular}{@{}lc@{}}
\toprule
\textbf{Parameter} & \textbf{Value} \\ \midrule
Particle count ($n$) & 2,660 \\
Smoothing radius ($h$) & 16.0 \\
Rest density ($\rho_0$) & 1.0 \\
Gas stiffness ($k$) & 2,000.0 \\
Viscosity ($\mu$) & 250.0 \\
Gravity ($g$) & 980.0 \\
Timestep ($\Delta t$) & 0.003 \\
\bottomrule
\end{tabular}
\end{table}

Each version was run until the simulation reached a steady state (particles settled into a puddle at the bottom). FPS was averaged across the entire session, and final particle positions were recorded to verify that the physics remained consistent across implementations.

\section{Results}

\subsection{Overall Performance}
Table~\ref{tab:results} shows the average FPS and speedup for each phase.

\begin{table}[h]
\centering
\caption{Benchmark Results (n=2660, Apple M1)}
\label{tab:results}
\begin{tabular}{@{}llcc@{}}
\toprule
\textbf{Phase} & \textbf{Config} & \textbf{Avg FPS} & \textbf{Speedup} \\ \midrule
P1: Baseline     & \texttt{-O0}, AoS, 1 thread     & 2.56   & 1.0$\times$  \\
P2: OpenMP       & \texttt{-O2}, AoS, 8 threads    & 58.89  & 23.0$\times$ \\
P3: NEON SIMD    & \texttt{-O2}, SoA, 1 thread     & 49.69  & 19.4$\times$ \\
P4: Combined     & \texttt{-O2}, SoA, 8 threads    & 57.83  & 22.6$\times$ \\ \bottomrule
\end{tabular}
\vspace{2mm}
\footnotesize{Note: FPS for Phases 2--4 is near the 60 FPS display cap, meaning actual compute throughput is higher than the reported frame rate suggests.}
\end{table}

A few things stand out. First, all three optimized versions hit close to 60 FPS, which is the display refresh cap---so FPS alone does not tell the full story. Second, the single-threaded SIMD version (49.69 FPS) gets about 85\% of the 8-thread OpenMP performance (58.89 FPS), which is a surprising result. Third, the combined version does not show a big FPS jump over OpenMP alone because both are already near the cap.

\subsection{Per-Loop Timing}
The more useful comparison is the per-loop physics timing, which is not capped by the display. Table~\ref{tab:timings} shows the breakdown.

\begin{table}[h]
\centering
\caption{Physics Loop Timing (per frame)}
\label{tab:timings}
\begin{tabular}{@{}lccc@{}}
\toprule
\textbf{Phase} & \textbf{Density (ms)} & \textbf{Force (ms)} & \textbf{Total (ms)} \\ \midrule
P2: OpenMP (AoS)          & 3.867  & 2.302  & 6.169  \\
P4: Combined (SoA + NEON) & 0.722  & 2.464  & 3.185  \\ \bottomrule
\end{tabular}
\end{table}

The combined version cuts the density loop time by 5.4$\times$ compared to OpenMP alone (3.867~ms $\to$ 0.722~ms). This is where the SoA layout and NEON vectorization make the biggest difference---the density kernel reads only position data, so the cache-friendly SoA layout and 4-wide vector loads are a perfect fit.

The force loop, however, barely improved (2.302~ms $\to$ 2.464~ms). This is because the force kernel reads density, pressure, and velocity in addition to position, and also requires square root and reciprocal operations---there is more data dependency and less opportunity for clean vectorization.

\subsection{Physics Validation}
To confirm that optimizations did not break the physics, we compared final particle positions across all four phases:

\begin{table}[h]
\centering
\caption{Final Particle Bounds (steady state)}
\label{tab:validation}
\begin{tabular}{@{}lccc@{}}
\toprule
\textbf{Phase} & \textbf{y min} & \textbf{y max} & \textbf{x range} \\ \midrule
Baseline  & 580.71 & 590.00 & 10--790 \\
OpenMP    & 581.13 & 590.00 & 10--790 \\
SIMD      & 580.88 & 590.00 & 10--790 \\
Combined  & 581.03 & 590.00 & 10--790 \\ \bottomrule
\end{tabular}
\end{table}

All versions converge to essentially the same final state---a thin puddle along the bottom boundary spanning the full width. The small differences ($< 0.5$ pixels) come from floating-point operation reordering and the fast reciprocal approximation in NEON.

\subsection{Baseline Scaling}
We also ran the baseline at two different particle counts to verify the $O(n^2)$ complexity:

\begin{table}[h]
\centering
\caption{Baseline Scaling with Particle Count}
\label{tab:scaling}
\begin{tabular}{@{}lcc@{}}
\toprule
\textbf{Particles} & \textbf{Avg FPS} & \textbf{Slowdown} \\ \midrule
1,500 & 6.5--6.7 & 1.0$\times$ \\
2,660 & 2.2--2.6 & $\sim$2.8$\times$ \\ \bottomrule
\end{tabular}
\end{table}

The theoretical $O(n^2)$ ratio is $(2660/1500)^2 = 3.14\times$. Our measured $\sim$2.8$\times$ is close---the gap is because rendering and integration are $O(n)$ and do not scale quadratically.

\section{Analysis}

\subsection{RQ1: OpenMP Speedup}
OpenMP achieved a 23$\times$ speedup over the \texttt{-O0} baseline, reaching 58.89 FPS with 8 threads. The density and force loops---the two $O(n^2)$ bottlenecks---are embarrassingly parallel on the outer loop because each particle's computation is independent.

It is important to note that this 23$\times$ number includes both the effect of multi-threading \emph{and} the switch from \texttt{-O0} to \texttt{-O2}. The \texttt{-O2} flag alone probably accounts for a significant portion of the speedup (through loop unrolling, better register allocation, and instruction scheduling). To isolate the pure threading contribution, we would need to run the single-threaded AoS version at \texttt{-O2} as well, which we plan to do in future work.

The M1's unified memory architecture helps here---all cores share the same physical memory without NUMA penalties, so thread scaling is relatively smooth. However, since the Icestorm cores run at 2.0 GHz vs.\ the Firestorm cores at 3.2 GHz, 8 threads do not give 8$\times$ scaling.

\subsection{RQ2: SIMD Performance}
Single-threaded NEON SIMD reached 49.69 FPS, which is about 85\% of the 8-thread OpenMP result. This is a striking finding: one carefully optimized core can nearly match 8 cores running unoptimized code.

The theoretical maximum SIMD speedup for 128-bit vectors on 32-bit floats is 4$\times$, but the actual gains come from multiple sources stacked together:
\begin{itemize}
    \item Processing 4 particles per inner loop iteration reduces loop overhead.
    \item Branchless masking (\texttt{vcltq\_f32} + \texttt{vandq\_u32}) eliminates branch mispredictions that stall the pipeline.
    \item The SoA layout enables full cache-line utilization for vector loads.
    \item The \texttt{-O2} compiler flag further optimizes the surrounding scalar code.
\end{itemize}

\subsection{RQ3: AoS vs.\ SoA and Cache Efficiency}
The memory layout change is arguably the most important optimization in this project. In AoS, each \texttt{Particle} struct is 32 bytes. A 64-byte cache line holds 2 particles. But the density kernel only needs position data (8 bytes per particle)---so 75\% of each cache line load is wasted on velocity, force, and other fields.

In SoA, the \texttt{pos\_x} array is contiguous. One cache line holds 16 consecutive x-positions. The density kernel reads \texttt{pos\_x[j..j+3]} and \texttt{pos\_y[j..j+3]}---two 16-byte loads that fully utilize the data brought into cache.

The per-loop timing data supports this. When we look at the density loop specifically---which reads only positions---the combined SoA+NEON version is 5.4$\times$ faster than the AoS OpenMP version (0.722~ms vs.\ 3.867~ms). The force loop, which accesses more fields and has less regular access patterns, shows almost no improvement.

The overall takeaway is that \textbf{how you lay out data in memory matters as much as, or more than, how you parallelize the computation.}

\subsection{RQ4: Architectural Bottlenecks}
The bottleneck shifts as we optimize:
\begin{enumerate}
    \item \textbf{Baseline:} Purely compute-bound. The \texttt{-O0} flag and scalar loops mean the CPU spends almost all its time on arithmetic for 7 million pair interactions per frame.
    \item \textbf{OpenMP:} Still somewhat compute-bound, but the density loop (3.87~ms) being 1.7$\times$ slower than the force loop (2.30~ms) suggests memory access patterns start to matter---the density kernel touches more data per unit of computation.
    \item \textbf{Combined:} With physics down to 3.185~ms per frame (theoretical max $\sim$313 FPS without display cap), we are approaching the memory bandwidth limit. Eight cores running SIMD all compete for the same $\sim$68 GB/s memory bus.
    \item \textbf{Fundamental limit:} Regardless of optimization, the $O(n^2)$ algorithm itself is the ceiling. At 2,660 particles we can hit 60 FPS, but scaling to 10,000+ particles would require a fundamentally different approach like spatial hashing to reduce complexity to $O(n \cdot k)$.
\end{enumerate}

\section{3D Extension}
As a supplementary experiment, we extended the simulation to 3D with an interactive rotatable cube container. The 3D version uses the same SoA + NEON optimization for physics.

An interesting finding here was that \emph{rendering}, not physics, became the bottleneck. Using Raylib's \texttt{DrawSphere} function (which draws a full 3D mesh per particle) gave only 18.26 FPS with 1,331 particles. After switching to a lightweight GL point rendering approach using \texttt{rlBegin(RL\_LINES)}, we achieved 58.27 FPS with 1,575 particles---a 3.2$\times$ improvement from the rendering change alone.

\section{Limitations}

There are several things we would do differently or add in a more thorough study:

\begin{itemize}
    \item \textbf{Compiler flag confound.} Our baseline uses \texttt{-O0} while the optimized versions use \texttt{-O2}. This means the reported speedup numbers include compiler optimization gains on top of the actual threading/SIMD gains. Running the baseline at \texttt{-O2} would give a cleaner comparison.
    \item \textbf{No hardware profiling.} We did not collect cache miss rates, branch misprediction counts, or memory bandwidth numbers from hardware counters. Tools like Apple Instruments or Linux \texttt{perf} would let us verify our cache efficiency claims with hard data.
    \item \textbf{Single particle count.} All optimized benchmarks use $n = 2660$. Testing at larger counts (5,000; 10,000; 50,000) would show where each optimization breaks down.
    \item \textbf{No thread scaling curve.} We only tested with 8 threads. A plot of speedup vs.\ thread count (1, 2, 4, 8) would reveal whether scaling is linear or starts to plateau.
\end{itemize}

\section{Future Work}

The most impactful next step would be implementing a \textbf{spatial hash grid} to reduce the algorithmic complexity from $O(n^2)$ to $O(n \cdot k)$, where $k$ is the average number of neighbors per particle. This would likely give a bigger speedup than all of our current optimizations combined, especially at higher particle counts.

Other planned work includes:
\begin{itemize}
    \item Collecting hardware performance counters to validate cache efficiency analysis.
    \item Testing scaling behavior up to 10,000+ particles.
    \item Exploring GPU acceleration via Apple's Metal API to move the computation off the CPU entirely.
    \item Generating benchmark comparison plots for a more visual presentation of the results.
\end{itemize}

\section{Conclusion}

In this project, we measured the impact of multi-threading and SIMD vectorization on a 2D SPH fluid simulation running on the Apple M1. Our main findings are:

\begin{enumerate}
    \item OpenMP multi-threading with 8 threads reaches 58.89 FPS (23$\times$ over the unoptimized baseline), bringing the physics computation to 6.169~ms per frame.
    \item Single-threaded NEON SIMD reaches 49.69 FPS (19.4$\times$), achieving about 85\% of the 8-thread OpenMP result on just one core.
    \item Combining both approaches reduces total physics time to 3.185~ms per frame, with the density loop specifically dropping 5.4$\times$ from the SoA restructuring.
    \item Memory layout (AoS vs.\ SoA) has at least as much impact as the choice of parallelization strategy. The SoA transformation was essential for making SIMD effective and also improved cache utilization significantly.
\end{enumerate}

The key lesson we took away from this project is that optimizing \emph{how} data is accessed matters as much as optimizing \emph{how much} computation is done in parallel. A single well-optimized core with cache-friendly data layout can come close to matching many cores running on poorly structured data.

\bibliographystyle{plain}
\bibliography{references}

\end{document}
